\section{TIẾN ĐỘ VÀ KẾT QUẢ GIỮA KỲ}\label{sec:tien_do}

Chương này báo cáo trạng thái từng hạng mục tại mốc \mocbaocao. Mỗi kết quả được xếp vào một trong ba mức: \textbf{Hoàn thành} (có bằng chứng kiểm chứng trên lab), \textbf{Đang thực hiện} (đã có một phần cấu hình hoặc kiểm thử) và \textbf{Chưa bắt đầu}. Số liệu nào chỉ đo một lần được ghi rõ là kết quả của một lần thử.

\subsection{Tổng hợp tiến độ theo hạng mục}

\begingroup
\scriptsize
\setlength{\tabcolsep}{3pt}
\begin{longtable}{|c|P{5.2cm}|P{2.2cm}|P{5.4cm}|}
    \caption{Trạng thái các hạng mục tại mốc giữa kỳ}\label{tab:tien_do} \\
    \hline
    \textbf{\#} & \textbf{Hạng mục} & \textbf{Trạng thái} & \textbf{Bằng chứng / ghi chú} \\
    \hline
    \endfirsthead
    \multicolumn{4}{c}{\textit{Bảng~\ref{tab:tien_do} (tiếp theo)}} \\
    \hline
    \textbf{\#} & \textbf{Hạng mục} & \textbf{Trạng thái} & \textbf{Bằng chứng / ghi chú} \\
    \hline
    \endhead
    \hline
    \endfoot
    1 & Khảo sát, xác định phạm vi và mục tiêu & Hoàn thành & Phạm vi 4 site + controller trên EVE-NG \\
    \hline
    2 & Lý thuyết SDN (OpenFlow, Ryu, OVS) & Hoàn thành & Đối chiếu với hành vi thực tế trong lab \\
    \hline
    3 & Lý thuyết SD-WAN (controller, OMP, BFD, TLOC) & Hoàn thành & Tài liệu giải thích tunnel SD-WAN \\
    \hline
    4 & Thiết kế kiến trúc 4 campus & Hoàn thành & Tài liệu thiết kế, Chương~\ref{sec:thiet_ke} \\
    \hline
    5 & Quy hoạch IP, VLAN, DHCP, ASN & Hoàn thành & Bảng~\ref{tab:vlan_campus}--\ref{tab:system_ip} \\
    \hline
    6 & Topology EVE-NG và quy trình đồng bộ cấu hình & Hoàn thành & 72 node, 105 network, 47 cấu hình nhúng; công cụ kiểm tra \texttt{.unl} báo ĐẠT; \texttt{configs/} đồng bộ theo running-config thật (10/10) \\
    \hline
    7 & Core/Dist/Access, VLAN, trunk, VRRP, OSPF & Hoàn thành & Firewall học đủ tuyến \texttt{10.1.x} qua OSPF \\
    \hline
    8 & Firewall ASAv HA và quản trị ASDM & Hoàn thành & \texttt{show failover}: Primary Active, Secondary Standby Ready (16/10); ASDM 7.20 mở được từ PC-Management (17/10). Chưa thử chuyển failover \\
    \hline
    9 & DHCP-Server và cấp phát end-to-end & Hoàn thành & 8/8 VPC Site 100 nhận địa chỉ qua relay (20/09); 7/7 VPC sau lần bật lại 16/10 \\
    \hline
    10 & Ryu + kết nối 6 OVS & Hoàn thành & 6/6 OVS kết nối controller; tắt ONOS (10/10) vì làm controller khựng \\
    \hline
    11 & Ứng dụng L2 SDN (học MAC theo VLAN, ACL) & Hoàn thành & Cây dữ liệu theo đồ thị liên kết, VLAN động, chính sách tập trung (04/10) \\
    \hline
    11b & Tự phục hồi SDN, fast-failover và Campus SDN Console & Hoàn thành & 6/6 kịch bản cắt link ĐẠT; fast-failover uplink Access (23/10): mất 1--2 gói thay vì 6--8~s; 33/33 unit test ĐẠT \\
    \hline
    12 & Ba controller SD-WAN và PKI & Hoàn thành & vManage, vSmart, vBond hoạt động, CA nội bộ \\
    \hline
    13 & Onboard 8 vEdge, BFD và control connection & Hoàn thành & 8/8 vEdge lên controller; BFD 72/72 up (12/10); thêm vEdge dự phòng ký sẵn tại Site 900 (08/10) \\
    \hline
    14 & Định tuyến liên site và chính sách SD-WAN & Hoàn thành & Service VPN 1, OMP; SLA, AAR, data policy, control policy \\
    \hline
    14b & Internet breakout và cảnh báo email & Hoàn thành & DIA tại chi nhánh, dự phòng qua MPLS (04/10); vManage gửi mail khi BFD/Control/OMP down (11/10) \\
    \hline
    15 & Chi nhánh 200/300/400: Brand-FW, switch, VPC & Hoàn thành & DHCP cục bộ, router-on-a-stick, ACL \texttt{SDWAN\_IN}, IP SLA theo dõi đường WAN \\
    \hline
    16 & vEdge chi nhánh join fabric, 2 hướng SP & Hoàn thành & 6 vEdge chi nhánh trong fabric \\
    \hline
    16b & Chi nhánh mới Site 500 (mạng phẳng) & Đang thực hiện (chậm) & Phía nhà cung cấp đã cấu hình BGP; vEdge-S500 chưa lên fabric, switch và máy trạm còn cấu hình mặc định \\
    \hline
    17 & Dịch vụ Web/Mail/DNS/Syslog & Hoàn thành phần chính & DNS, HTTPS, gửi/nhận mail từ chi nhánh qua SD-WAN \\
    \hline
    18 & Kiểm thử tích hợp end-to-end & Đang thực hiện & SDN đã đo đủ 5 mục tiêu (23/10); failover chi nhánh đo một lần mỗi site; còn đo lặp SD-WAN, failover FW HA, mất controller SD-WAN \\
    \hline
    19 & Báo cáo cuối kỳ & Đang thực hiện & Bản nháp các chương; báo cáo giữa kỳ này \\
    \hline
    20 & Slide và demo trực tiếp & Đang thực hiện & Đã có slide báo cáo tiến độ; slide bảo vệ và kịch bản demo dự kiến 11/2026 \\
\end{longtable}
\endgroup

Trong 23 hạng mục (gồm 3 hạng mục bổ sung 11b, 14b, 16b), 18 hạng mục đã hoàn thành, 1 hạng mục hoàn thành phần chính (dịch vụ) và 4 hạng mục đang thực hiện. So với mốc 01/10/2026, phần SDN đã chuyển từ ``chạy được'' sang ``có số liệu đo và tự phục hồi''. Khối lượng còn lại tập trung vào đo lặp phía SD-WAN, kiểm thử failover firewall, Site 500 và tài liệu.

\subsection{Kết quả đã đạt được}

\subsubsection{Hạ tầng Campus chính}

Core-SW1/2 cung cấp gateway VRRP cho các VLAN 10/20/30/40/90 và trao đổi tuyến OSPF với firewall active. Cặp ASAv hoạt động active--standby, cấu hình tự đồng bộ sang unit dự phòng; cả hai unit đều có địa chỉ quản trị trong VLAN 99 (\texttt{10.1.99.33} và \texttt{.34}). PC quản trị mở được ASDM 7.20 tới unit active sau khi đồng bộ giờ với firewall và cấp chứng thư tự ký đúng địa chỉ quản trị. DHCP Server (Windows Server 2012 R2) có bốn scope cho bốn VLAN người dùng; máy trạm Site 100 nhận đúng địa chỉ qua DHCP relay trên Core.

\subsubsection{SDN}

Bộ điều khiển Ryu tại \texttt{10.1.99.10} quản lý sáu OVS (Dist-SW1/2, Access-SW1--4) bằng OpenFlow 1.3. Từ đầu tháng 10, ứng dụng điều khiển được viết lại theo hướng tự phục hồi:

\begin{itemize}
    \item Cây chuyển tiếp dữ liệu được tính từ đồ thị liên kết do controller tự dò; khi một liên kết chết, controller tính lại cây.
    \item VLAN quản trị 99 được thiết kế lại để không phụ thuộc STP của Core; mỗi OVS có nhiều đường quản trị đồng thời.
    \item VLAN động và chính sách tập trung: thêm VLAN hoặc luật truy cập qua REST API của controller, áp xuống mọi switch.
    \item \textbf{Fast-failover trên data plane (23/10):} BFD chu kỳ 100~ms giữa các OVS phát hiện uplink chết; mỗi Access có một group OpenFlow \textit{fast-failover} gồm uplink chính và uplink dự phòng, nên switch tự chuyển đường mà không chờ controller.
    \item Campus SDN Console: giao diện web trên controller hiển thị sức khỏe kênh điều khiển, topology, và chạy tự động các kịch bản cắt link.
\end{itemize}

ONOS từng được cài song song làm giao diện giám sát nhưng đã tắt ngày 10/10 vì chiếm CPU làm Ryu khựng. Bảng~\ref{tab:sdn_ket_qua} tổng hợp kết quả đo ngày 23/10/2026.

\begin{table}[H]
    \centering
    \small
    \caption{Kết quả đo SDN Site 100 (23/10/2026)}
    \label{tab:sdn_ket_qua}
    \begin{tabularx}{\textwidth}{|P{3.6cm}|X|}
        \hline
        \textbf{Mục tiêu} & \textbf{Kết quả} \\
        \hline
        Hiệu năng chuyển mạch & 0\% mất gói; cùng VLAN khoảng 0,9~ms, khác VLAN trung vị 4~ms \\
        \hline
        Quản lý chính sách tập trung & Một luật áp xuống 4 Access trong 12--18~ms \\
        \hline
        Cấp phát VLAN mới & 9~ms qua controller, so với khoảng 36~s cho mỗi Core khi cấu hình tay \\
        \hline
        Tải controller & 8--15\% CPU \\
        \hline
        Mất một uplink Access (cắt thật) & Trước fast-failover gián đoạn 6--8~s; sau fast-failover mất 1--2 gói (khoảng 0,2~s) \\
        \hline
        Tắt hẳn Dist-SW1 & Gián đoạn 1,8~s (17/10, sau khi sửa thứ tự thông báo MAC); trước đó 10,2~s \\
        \hline
    \end{tabularx}
\end{table}

Hạn chế còn lại: khi hỏng hai liên kết cùng lúc hoặc mất liên kết Distribution--Core, lưu lượng vẫn phải chờ controller tính lại cây (khoảng 2--3,7~s). Controller khựng 3--4 lần mỗi phút do máy chủ EVE-NG dùng chung với nhiều người, nên các số đo có nhiễu.

\subsubsection{SD-WAN}

Fabric gồm ba controller và tám vEdge thuộc Site 100--400. Các kết quả đã kiểm chứng:

\begin{itemize}
    \item 8/8 vEdge có control connection tới vBond, vSmart, vManage; BFD 72/72 phiên up sau khi nâng vCPU cho vEdge (12/10).
    \item Service VPN 1 trên tám vEdge; OMP quảng bá subnet của bốn site; máy trạm các site ping thông nhau.
    \item Chính sách tập trung trên vSmart: Application-Aware Routing theo hai lớp SLA \texttt{SLA\_REALTIME} và \texttt{SLA\_BUSINESS}; data policy chặn Telnet/SSH từ chi nhánh vào Server Farm và DMZ; control policy ưu tiên vEdge1 của mỗi site. Counter của từng luật tăng đúng khi tạo lưu lượng thật.
    \item Failover chi nhánh: khi đường MPLS của vEdge1 chi nhánh đứt, Brand-FW phát hiện qua IP SLA và chuyển sang vEdge2. Thời gian khôi phục đo một lần mỗi site (26/09): Site 200 3,3~s, Site 300 4,4~s, Site 400 3,3~s; tự quay về khi đường chính sống lại.
    \item Internet breakout (04/10): chi nhánh ra Internet trực tiếp, dự phòng qua MPLS; VLAN quản trị, Server Farm và Site 900 bị giới hạn chiều ra.
    \item Cảnh báo email (11/10): vManage gửi mail khi BFD, control connection hoặc OMP down; đã thử bằng cách cắt vEdge2-S200.
    \item Một vEdge dự phòng tại Site 900 đã được ký chứng thư và đưa vào whitelist sẵn (08/10), lên fabric với 10 phiên BFD, dùng để mở chi nhánh mới nhanh.
\end{itemize}

Ngày 26/09/2026, nhóm thử AAR bằng cách tạo trễ nhân tạo 1000~ms trên một đường WAN của Site 100. Lưu lượng ICMP từ Site 200 đã chuyển khỏi đường hầm bị trễ. Do chu kỳ đo BFD được đặt dài (\texttt{poll-interval} 120~s), thời gian chuyển đường thực tế là vài phút; đây là kết quả của một lần thử.

Để giảm tải máy chủ, cụm SD-WAN được tắt khi không đo và bật lại theo quy trình đã kiểm chứng.

\subsubsection{Dịch vụ qua SD-WAN}

Từ máy trạm Linux tại Site 200, nhóm kiểm tra phân giải DNS bản ghi \texttt{mail}/\texttt{MX}, truy cập HTTPS tới Web Server trong DMZ (máy chủ thấy địa chỉ thật của máy trạm, không bị NAT) và gửi/nhận thư qua cổng 587/IMAP. Firewall Site 100 chỉ mở đúng các dịch vụ cần thiết từ SD-WAN vào DMZ; các cổng khác bị chặn và ghi log.

\subsubsection{Công cụ đo kiểm}

Nhóm xây dựng một ứng dụng web chạy trên máy Windows tại Site 900. Ứng dụng chỉ lấy dữ liệu qua REST API của vManage và ping trực tiếp, không đọc qua máy chủ ảo hóa, để giống với cách vận hành thật. Các tab gồm: Tổng quan, Đo lường (độ trễ, jitter, mất gói giữa các site), SD-WAN Failover (đo thời gian khôi phục khi cắt link WAN), Độ tin cậy (tổng hợp RTO, tỉ lệ gói đến đích, thời gian khôi phục controller, độ sẵn sàng), Thông lượng, Kiểm chứng VLAN, Chính sách tập trung và Ký chứng thư vEdge. Tab Chính sách tập trung đã chạy trên lab thật. Ngày 24/10, lần chạy thật đầu tiên của tab SD-WAN Failover phát hiện lỗi khi gọi API ping của vManage; lỗi đã được sửa và đang chờ đo lại.

\subsection{Vấn đề kỹ thuật đã gặp và cách xử lý}

Bảng~\ref{tab:van_de} tổng hợp các sự cố tiêu biểu. Mỗi sự cố được xác định nguyên nhân bằng chẩn đoán theo lớp trước khi sửa.

\begingroup
\scriptsize
\setlength{\tabcolsep}{3pt}
\begin{longtable}{|P{3.3cm}|P{4.7cm}|P{4.8cm}|}
    \caption{Một số vấn đề kỹ thuật và cách xử lý}\label{tab:van_de} \\
    \hline
    \textbf{Hiện tượng} & \textbf{Nguyên nhân} & \textbf{Cách xử lý} \\
    \hline
    \endfirsthead
    \multicolumn{3}{c}{\textit{Bảng~\ref{tab:van_de} (tiếp theo)}} \\
    \hline
    \textbf{Hiện tượng} & \textbf{Nguyên nhân} & \textbf{Cách xử lý} \\
    \hline
    \endhead
    \hline
    \endfoot
    Máy trạm Site 100 không nhận DHCP qua OVS & Luật table-miss không gửi gói lên controller; ứng dụng đọc VLAN sai trường (\texttt{dl\_vlan} của OpenFlow 1.0 thay vì \texttt{vlan\_vid}) & Sửa luật table-miss và cách so khớp VLAN trong ứng dụng Ryu \\
    \hline
    Bão broadcast mỗi khi một OVS khởi động & OVS tự cài flow \textit{in-band} ẩn với action NORMAL khi chưa nối được controller, làm ARP bị flood & Tắt in-band control (\texttt{disable-in-band=true}) trên cả 6 OVS \\
    \hline
    Switch mất kết nối controller & VLAN 99 có vòng lặp khi tắt STP trên OVS & Thiết kế lại VLAN 99 thành cây không vòng lặp, không phụ thuộc STP của Core \\
    \hline
    Mất một uplink Access gián đoạn 6--8~s & Lưu lượng chỉ chuyển đường khi controller phát hiện liên kết chết và tính lại cây & BFD giữa OVS + group OpenFlow fast-failover: switch tự chuyển uplink \\
    \hline
    Tắt Dist-SW1 gián đoạn 10~s & Gói thông báo MAC phát trước khi Dist-SW2 gỡ luật chặn cổng xuống Access, Core giữ MAC cũ & Chỉ phát thông báo sau khi switch xác nhận đã cài xong luật; còn 1,8~s \\
    \hline
    Controller Ryu khựng nhiều giây & ONOS chạy cùng máy chiếm CPU; máy chủ EVE-NG dùng chung tải cao & Tắt ONOS, tăng ưu tiên CPU cho Ryu; thêm thời gian ân hạn trước khi kết luận liên kết chết \\
    \hline
    vEdge kẹt ở trạng thái \textit{connect} & Chứng thư có trường O khác ``Cisco Systems'' nên vBond/vSmart từ chối & Ký lại chứng thư với subject đúng chuẩn \\
    \hline
    vEdge rơi khỏi fabric sau khi controller khởi động lại & Danh sách thiết bị được phép trên vBond chỉ lưu trong bộ nhớ & Đẩy lại danh sách từ vManage; kiểm tra whitelist sau mỗi lần khởi động lại controller \\
    \hline
    Màu \texttt{biz-internet} không có đường tới controller & Route /24 qua MPLS che route /16 qua Internet (longest-prefix) & Thêm route /24 cùng độ dài qua cả hai transport \\
    \hline
    Tunnel Internet--MPLS down hàng loạt & Khai sai màu TLOC trên vEdge2 chi nhánh & Sửa màu đúng với loại đường truyền \\
    \hline
    BFD chéo màu Site 100 down & Mỗi TLOC chỉ có đường tới transport còn lại qua một cổng & Thêm default route cho từng transport trong VPN 0 \\
    \hline
    Failover chi nhánh mất 100\% gói & Gói về đi vào vEdge khác chiều đi, firewall ASA hủy do bất đối xứng; Brand-FW không biết đường qua vEdge1 đã chết & Control policy ưu tiên vEdge1, sửa route trên vEdge2, IP SLA + track trên Brand-FW \\
    \hline
    Cổng WAN của vEdge không nhận gói (gặp nhiều lần) & Card mạng ảo bị treo phía nhận & \texttt{shutdown}/\texttt{no shutdown} cổng; kiểm tra bộ đếm drop sau mỗi lần khởi động vEdge \\
    \hline
    Mail ``BFD Node Down'' ngắn vài giây & vEdge chỉ có 1 vCPU, CPU quá tải làm BFD lỡ chu kỳ & Nâng vEdge lên 2 vCPU; theo dõi lại sau khi hạ một số vEdge về 1 vCPU để giảm tải máy chủ \\
    \hline
    ASDM báo chứng thư hết hạn và sai hostname & Giờ PC quản trị lệch múi giờ so với firewall; chứng thư tự sinh không chứa địa chỉ quản trị & Đồng bộ giờ, tạo chứng thư tự ký theo địa chỉ \texttt{10.1.99.33} \\
\end{longtable}
\endgroup

\subsection{Đánh giá giữa kỳ}

Tiến độ thực tế bám sát kế hoạch: khảo sát, thiết kế và triển khai cốt lõi đã hoàn thành; phần SDN đã có bộ số liệu đo đủ năm mục tiêu và cơ chế tự phục hồi được kiểm chứng bằng cắt link thật. Mọi kết quả đều có bằng chứng từ thiết bị và các sự cố đều được ghi nhận nguyên nhân gốc, tạo thành tài liệu vận hành có giá trị.

Các hạn chế hiện tại:

\begin{itemize}
    \item Phía SD-WAN chưa có số liệu đo lặp lại: thời gian failover mới đo một lần mỗi site; thời gian khôi phục controller và độ sẵn sàng chưa đo.
    \item Chưa kiểm thử chuyển failover của cặp firewall ASAv.
    \item SDN chỉ có một controller Ryu, là điểm lỗi đơn của mặt phẳng điều khiển; hỏng kép vẫn phụ thuộc controller.
    \item Hai đợt toàn bộ vEdge mất kết nối (đêm 22--23/09 và rạng sáng 12/10) chưa xác định được nguyên nhân; lần sau sẽ thu bằng chứng trước khi khởi động lại.
    \item Site 500 chậm so với kế hoạch (hạn ban đầu 10/10).
    \item Máy chủ EVE-NG dùng chung nên không chạy đồng thời toàn bộ 72 node; số đo chịu nhiễu khi máy chủ tải cao.
\end{itemize}
